%=======================================================================
%  2103304 Automatic Control I -- Course Syllabus (Fall 2026)
%  Department of Mechanical Engineering, Chulalongkorn University
%  Compiles with pdfLaTeX.
%=======================================================================
\documentclass[11pt]{article}

\usepackage[T1]{fontenc}
\usepackage[utf8]{inputenc}
\usepackage[margin=1in]{geometry}
\usepackage{array}
\usepackage{booktabs}
\usepackage{tabularx}
\usepackage{enumitem}
\usepackage{xcolor}
\usepackage{colortbl}
\usepackage{titlesec}
\usepackage{mathptmx}      % Times-like serif, clean academic look
\usepackage{fancyhdr}
\usepackage[hidelinks]{hyperref}

% ---- Document version / date (edit here) --------------------------
\newcommand{\docversion}{1.0}
\newcommand{\docdate}{24 July 2026}

% ---- Brand colors (Chula ME) --------------------------------------
\definecolor{brandred}{HTML}{8B2332}
\definecolor{brandred20}{HTML}{ECD7D9}
\definecolor{brandgray}{HTML}{6F6D6E}

% ---- Section styling ----------------------------------------------
\titleformat{\section}
  {\color{brandred}\normalfont\Large\bfseries}{}{0pt}{}[{\color{brandred!40}\titlerule[1pt]}]
\titlespacing*{\section}{0pt}{14pt}{6pt}

\titleformat{\subsection}
  {\color{brandred}\normalfont\large\bfseries}{}{0pt}{}
\titlespacing*{\subsection}{0pt}{10pt}{3pt}

\setlist[enumerate]{leftmargin=1.5em, itemsep=3pt, topsep=3pt}
\setlist[itemize]{leftmargin=1.4em, itemsep=2pt, topsep=3pt}

\newcolumntype{L}[1]{>{\raggedright\arraybackslash}p{#1}}
\newcolumntype{C}[1]{>{\centering\arraybackslash}p{#1}}

\renewcommand{\arraystretch}{1.25}
\setlength{\parindent}{0pt}
\setlength{\parskip}{4pt}

% Prevent awkward hyphenation of these words (no break points listed)
\hyphenation{mechanical mechanics mechanical-engineering}
% Absorb occasional tight lines by micro-stretching instead of overflowing
\setlength{\emergencystretch}{1em}

% PI tag used to map CLOs to program performance indicators
\newcommand{\pitag}[1]{\hfill{\small\color{brandred}\textbf{[#1]}}}

% Light horizontal rule between weeks in the schedule
\newcommand{\weekline}{\arrayrulecolor{black!18}\hline\arrayrulecolor{black}}

% ---- Running header / footer on every page ------------------------
\pagestyle{fancy}
\fancyhf{}
\renewcommand{\headrulewidth}{0.4pt}
\renewcommand{\footrulewidth}{0pt}
\fancyhead[L]{\small\color{brandgray}2103304 Automatic Control I: Syllabus (Fall 2026)}
\fancyhead[R]{\small\color{brandgray}Chawin Ophaswongse, Ph.D.}
\fancyfoot[L]{\small\color{brandgray}Version \docversion}
\fancyfoot[C]{\small\color{brandgray}\thepage}
\fancyfoot[R]{\small\color{brandgray}\docdate}
\setlength{\headheight}{14pt}
\setlength{\headsep}{18pt}

\begin{document}

%=======================================================================
%  Title block
%=======================================================================
\begin{center}
    {\LARGE\bfseries Course Syllabus}\\[6pt]
    {\large Department of Mechanical Engineering}\\
    {\large Chulalongkorn University}\\[4pt]
    {\normalsize Academic Year: Fall 2026 (First Semester)}
\end{center}
% \vspace{4pt}
% {\color{brandred!50}\hrule height 1.5pt}
% \vspace{10pt}

%=======================================================================
\section*{\color{brandred}Course Information}
\vspace{-2pt}
\begin{tabular}{@{}>{\bfseries}L{5.0cm} L{11.0cm}@{}}
    Course Number   & 2103304 \\
    Course Title    & Automatic Control I \\
    Credits         & 3 (3--0--9) \\
    Semester        & Fall 2026 (First Semester) \\
    Instructor      & Chawin Ophaswongse, Ph.D. \\
    Co-Instructor   & Assoc.\ Prof.\ Thitima Jintanawan, Ph.D. \\
    Teaching Assistants (TAs) & \dots\\
    Class Time      & Monday, 13:00--16:00 \\
    Location        & Room 409, Building 3, Faculty of Engineering \\
    Office Hours    & TBD, Hans Bunthli Building, 4th Floor, Room 407 \\
    Prerequisite    & \dots \\
\end{tabular}

%=======================================================================
\section{Course Description}
This course provides a foundational understanding of feedback control systems with
applications in mechanical engineering. It develops the full analysis-and-design
workflow: representing a physical system, deriving its governing equations, and
building the engineering models needed to predict and shape its behavior. Core
representations include ordinary differential equations, transfer functions, and
state-space models, developed from mechanical and electro-mechanical examples such
as mass--spring--damper systems and DC-motor drives. Models are obtained both from
first principles and experimentally through \emph{system identification}, fitting a
model to measured input--output data.

Students learn to analyze system behavior in both the time and frequency domains
(transient response, steady-state error, and stability) and to design and tune
controllers (P, PI, PD, PID, lead/lag, and full-state feedback) to meet performance
specifications such as rise time, overshoot, and settling time. Classical techniques
including block-diagram reduction, the Routh--Hurwitz criterion, root locus, and
Bode/Nyquist frequency response are introduced alongside an entry-level treatment of
state-space design.

Analysis and design are practiced throughout using \textbf{MATLAB/Simulink} with the
Control System Toolbox. The \emph{Servo Mini-labs} and a \emph{Final Project} connect
theory to hardware through a MATLAB script or Simulink-based control interface to a real servo plant
(a brushed DC-motor system). The emphasis is on control \emph{analysis and design}:
identifying the plant model from experiment (system identification), designing and
tuning the controller, and interpreting the measured response, rather than low-level
firmware or embedded programming, so students focus on control-engineering decisions
instead of implementation plumbing.

%=======================================================================
\newpage
\section{Course Learning Outcomes (CLOs)}
On successful completion of this course, students will be able to:

\begin{center}
\begin{tabularx}{\textwidth}{@{}C{1cm} X C{1.2cm}@{}}
\toprule
\textbf{CLO} & \textbf{Course Learning Outcome} & \textbf{PI} \\
\midrule
1 & Identify the relevant parameters and derive the governing equations of a mechanical or electro-mechanical system, and lay out a systematic procedure for analyzing a real engineering problem (equations of motion, transfer functions, free/damped vibration, stability). & 1.2 \\
2 & Construct engineering models of a given problem or real system suitable for analysis (free-body diagrams, kinematics/kinetics diagrams, control block diagrams) and translate them into differential-equation, transfer-function, and state-space representations. & 1.3 \\
3 & Analyze system behavior in the time and frequency domains, including transient response, steady-state error, and stability. & 1.2 \\
4 & Design and tune feedback controllers (P/PI/PD/PID, lead/lag, full-state feedback) to meet performance specifications, using root locus and frequency-response methods. & 1.2 \\
5 & Apply MATLAB/Simulink to identify, model, simulate, analyze, and design control systems, and validate a controller on a real servo plant through a Simulink hardware interface. & 1.3 \\
\bottomrule
\end{tabularx}
\end{center}

\subsection{Program Performance Indicators (PIs) Addressed}
\begin{itemize}
    \item \textbf{PI 1.2 (Parameters and governing equations in applied mechanics).}
          Ability to identify variables, write the governing (control) equations, and
          define the procedure for analyzing engineering problems and real mechanical
          systems (e.g., static equilibrium, stress--strain in beams, stress
          transformation, failure theory, equation of motion, transfer function,
          free/damped vibration, and stability).
    % TH: สามารถระบุตัวแปร เขียนสมการควบคุม พร้อมทั้งกำหนดขั้นตอนในการวิเคราะห์โจทย์และระบบจริงด้านกลศาสตร์
    \item \textbf{PI 1.3 (Engineering modeling).} Ability to build mechanical-engineering
          models of problems and real systems for use in analysis (e.g., free-body
          diagram, kinematics diagram, kinetics diagram, stress--strain element, control
          block diagram, P--V diagram, flow diagram, and energy-transfer diagram).
    % TH: สามารถสร้างแบบจำลองทางวิศวกรรมเครื่องกลของโจทย์ปัญหาและระบบจริงเพื่อนำไปใช้ในการวิเคราะห์ปัญหาได้
\end{itemize}

%=======================================================================
\section{Textbooks}
\textbf{Primary}
\begin{itemize}
    \item Nise, N. S. (2011). \emph{Control Systems Engineering} (6th ed.). Wiley.
          \\\hfill\textit{(Beginner-friendly)}
\end{itemize}
\textbf{Supplementary}
\begin{itemize}
    \item Ogata, K. (2010). \emph{Modern Control Engineering} (5th ed.). Prentice Hall.
          \\\hfill\textit{(State space; MATLAB; more mathematically rigorous)}
    \item Dorf, R. C., \& Bishop, R. H. (2017). \emph{Modern Control Systems} (13th ed.). Pearson.
          \\\hfill\textit{(Design emphasis; rich in illustrative examples)}
\end{itemize}

%=======================================================================
\newpage
\section{Course Schedule}
\begin{center}
\scriptsize
\setlength{\tabcolsep}{4pt}
\begin{tabularx}{\textwidth}{@{}C{0.6cm} C{1.3cm} L{2.9cm} X C{0.7cm} L{3.7cm}@{}}
\toprule
\textbf{Wk} & \textbf{Date} & \textbf{Topic} & \textbf{Key Content} & \textbf{HW} & \textbf{Servo Mini-lab} \\
\midrule
1  & 3 Aug  & Intro \& Math Primer   & Open-/closed-loop $\cdot$ ODEs $\cdot$ Laplace     & HW1 & \\
\weekline
2  & 10 Aug & System Modeling I      & LTI $\cdot$ transfer functions $\cdot$ elec./mech.  & & Lab 1: Intro \& Setup \\
\weekline
3  & 17 Aug & System Modeling II     & State space $\cdot$ nonlinearities $\cdot$ linearization & HW2 & \\
\weekline
4  & 24 Aug & Time Response          & 1st/2nd-order $\cdot$ transient-response specs     & HW3 & Lab 2: Time Response \\
\weekline
5  & 31 Aug & Block Diagrams         & Reduction $\cdot$ equivalent forms $\cdot$ feedback & & \\
\weekline
6  & 7 Sep  & Stability              & Routh--Hurwitz $\cdot$ gain range for stability   & HW4 & Lab 3: P-Control \& Stability \\
\weekline
7  & 14 Sep & Midterm Review         & Review of Weeks 1--6                               & & \\
\midrule
\multicolumn{6}{@{}l}{\textit{Midterm Examination: Wednesday, 23 September 2026, 13:00--16:00}} \\
\midrule
8  & 28 Sep & Steady-State Errors    & System type $\cdot$ gain design $\cdot$ disturbances & HW5 & \\
\weekline
9  & 5 Oct  & Root Locus I           & Sketching rules $\cdot$ transient response via gain & HW6 & \\
\weekline
10 & 12 Oct & Root Locus II          & PI/PD/PID $\cdot$ lead/lag compensation            & & Lab 4: Root Locus \\
\weekline
11 & 19 Oct & Frequency Response I    & Asymptotic Bode plots $\cdot$ bandwidth           & HW7 & \\
\weekline
12 & 26 Oct & Frequency Response II   & Nyquist \& margins $\cdot$ lead/lag $\cdot$ PID     & & Lab 5: Frequency Response \\
\weekline
13 & 2 Nov  & Design via State Space & Controllability $\cdot$ full-state feedback         & HW8 & Lab 6: Pole Placement \\
\weekline
14 & 9 Nov  & Project Demo Day       & Group presentations $\cdot$ control competition   & & \\
\weekline
15 & 16 Nov & Final Review           & Review of all topics                              & & \\
\midrule
\multicolumn{6}{@{}l}{\textit{Final Examination: Monday, 30 November 2026, 08:30--11:30}} \\
\bottomrule
\end{tabularx}
\end{center}

%=======================================================================
\section{Course Policies}
\begin{itemize}
    \item \textbf{Communication.} Announcements and materials are posted on MyCourseVille;
          day-to-day discussion is on the course Discord.
    \item \textbf{Attendance.} Mini-lab sessions are mandatory and graded.
    \item \textbf{Mini-labs.} Pre-lab handouts are released one week before each mini-lab
          (reading required before the session). Lab reports (groups of up to 4) are due
          one week after the lab session.
    \item \textbf{Assignments.} Homework is due 10 days after it is assigned unless
          otherwise specified. Late work is penalized per the stated policy.
    \item \textbf{Academic integrity.} Collaboration is encouraged, but submitted work
          must be each group's own; plagiarism and exam misconduct follow Faculty of
          Engineering regulations.
    \item \textbf{Software.} Install MATLAB/Simulink and the Control System Toolbox
          before Week 2.
\end{itemize}


%=======================================================================
\newpage
\section{Assessment}
\begin{center}
\begin{tabular}{@{}L{9cm} C{2cm}@{}}
\toprule
\textbf{Component} & \textbf{Weight} \\
\midrule
Assignments (Homework)      & 10\% \\
Servo Mini-labs             & 10\% \\
Midterm Examination         & 30\% \\
Final Project               & 15\% \\
Final Examination           & 35\% \\
\midrule
\textbf{Total}              & \textbf{100\%} \\
\bottomrule
\end{tabular}
\end{center}

%=======================================================================
\section{Servo Mini-labs}
Six 45-minute hands-on sessions run in parallel with the lectures. Students identify a
plant model from experiment (system identification), model its components (actuator,
sensor, controller), then design and tune controllers in \textbf{MATLAB/Simulink} and
validate them on a real-time servo setup: a brushed DC-motor with an interchangeable
inertia-disk or pendulum load. Interaction with the hardware is entirely through the
Simulink control interface (\emph{no low-level or embedded programming is required}),
so the focus stays on control analysis and design. Sessions are mandatory and graded;
a laptop with MATLAB/Simulink and the Control System Toolbox is required.

\textbf{No-load vs.\ inertia-disk comparison (Labs 2--6).} The inertia disk is used from
Lab 2 onward as a single-parameter change to the plant: mounting it raises the inertia
$J$, which scales the mechanical time constant $\tau = J/(b + K_tK_b/R)$ while leaving
the DC gain unchanged. Each session therefore begins with a short no-load baseline test
on the bare shaft, repeated with the disk attached, and every lab report includes an
overlay of the two responses with the resulting parameter change quantified.

\emph{Tentative session plan (subject to adjustment):}
\begin{itemize}
    \item \textbf{Lab 1: Intro \& Setup (Wk 2).} Connect the DC-motor/inertia-disk
          plant to Simulink; get familiar with the control interface, encoder sensing,
          and actuation; run open-loop step tests and a first system-identification of
          the motor parameters (no-load baseline).
    \item \textbf{Lab 2: Time Response (Wk 4).} Measure open- and closed-loop step
          responses; extract first-/second-order parameters and transient specifications
          (rise time, overshoot, settling time) and compare against the identified model.
          \emph{No-load vs.\ disk:} the steady-state speed is unchanged while $\tau$
          increases; cross-check the identified inertia against the geometric value
          $J_d = \frac{1}{2}mr^2$ and the predicted
          $\tau_{\mathrm{disk}} = \tau_0(1 + J_d/J_0)$.
    \item \textbf{Lab 3: P-Control \& Stability (Wk 6).} Implement proportional
          position/speed control; observe how gain shapes transient response and the
          onset of instability. \emph{No-load vs.\ disk:} hold $K_p$ fixed across both
          conditions and explain the increased overshoot and ringing from
          $\omega_n \propto 1/\sqrt{J}$ and $\zeta \propto 1/\sqrt{J}$.
    \item \textbf{Lab 4: Root Locus (Wk 10).} Design controller gain (P/PD/PID) via
          root locus to meet transient specifications and validate on the plant.
          \emph{No-load vs.\ disk:} apply the gain designed for the loaded plant to the
          unloaded plant and observe the specification violation (motivating robustness).
    \item \textbf{Lab 5: Frequency Response (Wk 12).} Measure the plant's frequency
          response (Bode); read gain/phase margins and tune a lead/lag or PID controller.
          \emph{No-load vs.\ disk:} compare bandwidth and the gain/phase margins of a
          fixed controller under both load conditions.
    \item \textbf{Lab 6: Pole Placement (Wk 13).} Design full-state feedback by pole
          placement for the DC-motor with the pendulum load and validate the closed-loop
          response. \emph{No-load vs.\ disk:} design for each load condition separately
          and show that feedback can render the two closed-loop responses nearly
          identical, subject to the available actuator authority.
\end{itemize}

%=======================================================================
\section{Final Project}
Working in groups of up to four (beginning after the midterm), students carry out the
complete workflow on a given electro-mechanical plant: performing system identification
to obtain the plant model, building a \textbf{Simulink simulation model} and validating
it against the measured hardware response, then designing, tuning, and demonstrating a
controller to meet target performance specifications through the MATLAB/Simulink
hardware interface. Comparing simulation against hardware is a central theme: students
are expected to explain and reconcile the differences they observe.

\emph{Tentative milestones (subject to adjustment):}
\begin{itemize}
    \item \textbf{Milestone 1: System ID \& Simulation Model (Week 9, 5 Oct).} Run
          experiments to identify the plant parameters; build a Simulink simulation model
          and validate it against the measured open-loop hardware response.
    \item \textbf{Milestone 2: Controller Design in Simulation (Week 11, 19 Oct).}
          Specify target performance, design and tune the controller in simulation, and
          verify that it meets the specifications on the simulation model.
    \item \textbf{Milestone 3: Hardware Implementation \& Tuning (Week 13, 2 Nov).}
          Deploy the controller on the real plant, retune as needed, and document the
          discrepancies between simulation and hardware.
    \item \textbf{Final Deliverable: Demo Day \& Report (Week 14, 9 Nov).} Present the
          final closed-loop performance in a demo-day competition and submit the final
          project report.
\end{itemize}

\vspace{6pt}
{\small\color{brandgray}\textit{This syllabus is subject to minor adjustment;
any changes will be announced in class and on MyCourseVille.}}


\end{document}
